\subsection{代数}

设 \(\mathbf{M}\) 是 \(\mathbf{K}\) 上的含单位元模，并带有一个双线性映射 \((x,y)\mapsto xy\)，它从 \(\mathbf{M}\times \mathbf{M}\) 到 \(\mathbf{M}\)。除乘法的结合律外，代数的所有公理都满足。滥用术语，称 \(\mathbf{M}\) 为 \(\mathbf{K}\) 上不必结合的代数；有时在不会引起混淆时，也简称为 \(\mathbf{K}\) 上的代数。本号中采用后一种记法。

若在 K-模 M 上赋予乘法 \((x,y)\mapsto yx\)，便得到一个代数，称为上述代数的逆代数。

M 的一个在乘法下稳定的子 K-模 N 以显然的方式成为 K 上的代数。称 N 为 M 的子代数。如果 \(x \in N\)、\(y \in M\) 蕴含 \(yx \in N\)（分别为 \(xy \in N\)），则称 N 为 M 的左（分别为右）理想。如果 N 既是左理想又是右理想，则称 N 为 M 的双边理想。在这种情形，利用 M 上的乘法，取商后可在商模 M/N 上定义双线性乘法，使 M/N 具有代数结构。称 M/N 为 M 关于 N 的商代数。

设 \(\mathbf{M}_1\) 和 \(\mathbf{M}_2\) 是 \(\mathbf{K}\) 上的两个代数，\(\phi\) 是从 \(\mathbf{M}_1\) 到 \(\mathbf{M}_2\) 的映射。称 \(\phi\) 为同态，当且仅当 \(\phi\) 是 K-线性的，并且 \(\phi(xy) = \phi(x)\phi(y)\) 对 \(x\in \mathbf{M}_1\)、\(y\in \mathbf{M}_1\) 成立。\(\mathbf{N}\) 是 \(\phi\) 的核，且是 \(\mathbf{M}_1\) 的双边理想；\(\phi\) 的像是 \(\mathbf{M}_2\) 的子代数。取商后，\(\phi\) 定义了从代数 \(\mathbf{M}_1 / N\) 到代数 \(\phi (\mathbf{M}_1)\) 的同构。

设 \(\mathbf{M}\) 是 \(\mathbf{K}\) 上的代数。映射 \(\mathbf{D}\) 从 \(\mathbf{M}\) 到 \(\mathbf{M}\)，称为 \(\mathbf{M}\) 的导子，若它是 \(\mathbf{K}\)-线性的，并且 \(\mathrm{D}(xy) = (\mathrm{D}x)y + x(\mathrm{D}y)\) 对所有 \(x \in \mathbf{M}\) 和 \(y \in \mathbf{M}\) 成立。这个定义推广了 Algebra, Chapter IV, § 4, no. 3 中的定义 3。\(\mathbf{M}\) 的导子的核是 \(\mathbf{M}\) 的子代数。如果 \(\mathrm{D}_1\) 和 \(\mathrm{D}_2\) 是 \(\mathbf{M}\) 的导子，则 \(\mathrm{D}_1\mathrm{D}_2 - \mathrm{D}_2\mathrm{D}_1\) 也是 \(\mathbf{M}\) 的导子（参见 Algebra, Chapter IV, § 4, no. 3, Proposition 5；该命题的证明不使用代数的结合律）。

设 \(\mathbf{M}_1\) 和 \(\mathbf{M}_2\) 是 K 上的两个代数。在积 K-模 \(\mathbf{M} = \mathbf{M}_1 \times \mathbf{M}_2\) 上规定乘法如下：

\[
(x _ {1}, x _ {2}) (y _ {1}, y _ {2}) = (x _ {1} y _ {1}, x _ {2} y _ {2}),
\]

对所有满足 \(x_1, y_1\) 属于 \(\mathbf{M}_1, x_2, y_2\) 属于 \(\mathbf{M}_2\) 的情形，如此定义的代数称为 \(\mathbf{M}_1\) 和 \(\mathbf{M}_2\) 的积代数。映射 \(x_1 \mapsto (x_1, 0)\)（分别为 \(x_2 \mapsto (0, x_2)\)）将 \(\mathbf{M}_1\)（分别为 \(\mathbf{M}_2\)）同构到 M 的一个双边理想上。在这些同构下，把 \(\mathbf{M}_1\) 和 \(\mathbf{M}_2\) 视为 M 的双边理想。于是 K-模 M 是 \(\mathbf{M}_1\) 与 \(\mathbf{M}_2\) 的直和。反之，设 M 是 K 上的代数，\(\mathbf{M}_1, \mathbf{M}_2\) 是 M 的两个双边理想，且 M 是 \(\mathbf{M}_1\) 与 \(\mathbf{M}_2\) 的直和。则 \(\mathbf{M}_1\mathbf{M}_2 \subset \mathbf{M}_1 \cap \mathbf{M}_2 = \{0\}\)；于是若 \(x_1, y_1\) 属于 \(\mathbf{M}_1\)，\(x_2, y_2\) 属于 \(\mathbf{M}_2\)，则 \((x_1 + x_2)(y_1 + y_2) = x_1y_1 + x_2y_2\)，从而可将 M 视为积代数 \(\mathbf{M}_1 \times \mathbf{M}_2\)。\(\mathbf{M}_1\) 的每个左（分别为右、双边）理想都是 M 的左（分别为右、双边）理想。任意有限族代数的情形中相应结果如何表述，留给读者完成。

设 \(\mathbf{M}\) 是 \(\mathbf{K}\) 上的代数，并设 \(\mathbf{K}\)-模 \(\mathbf{M}\) 有一个基 \((a_{\lambda})_{\lambda \in \mathbf{L}}\)。存在唯一的系统 \((\gamma_{\lambda \mu \nu})_{(\lambda, \mu, \nu) \in \mathbf{L} \times \mathbf{L} \times \mathbf{L}}\)，其元素属于 \(\mathbf{K}\)，并满足 \(a_{\lambda}a_{\mu} = \sum_{\nu} \gamma_{\lambda \mu \nu}a_{\nu}\) 对 \(\lambda, \mu\) 属于 \(\mathbf{L}\) 成立。称 \(\gamma_{\lambda \mu \nu}\) 为 \(\mathbf{M}\) 关于基 \((a_{\lambda})\) 的结构常数。

设 \(\mathbf{M}\) 是 \(\mathbf{K}\) 上的代数，\(\mathbf{K}_0\) 是含单位元的交换环，\(\rho\) 是将单位元映到单位元的从 \(\mathbf{K}_0\) 到 \(\mathbf{K}\) 的同态。则可将 \(\mathbf{M}\) 视为 \(\mathbf{K}_0\) 上的代数，规定 \(\alpha.x = \rho(\alpha).x\)（其中 \(\alpha \in \mathbf{K}_0, x \in \mathbf{M}\)）。特别地，当 \(\mathbf{K}_0\) 是含单位元的 \(\mathbf{K}\) 的子环，并取 \(\rho\) 为 \(\mathbf{K}_0\) 到 \(\mathbf{K}\) 的包含映射时，正是这种情形。

设 M 是 K 上的代数，\(K_{1}\) 是含单位元的交换环，\(\sigma\) 是将单位元映到单位元的从 K 到 \(K_{1}\) 的同态。令 \(M_{(K_{1},\sigma)} = M_{(K_{1})}\) 表示 \(K_{1}\)-模，即由 M 将标量环扩张到 \(\mathbf{K}_1\) 后得到的模（参见 Algebra, Chapter II, § 5）。M 上的乘法自然定义了一个 \(\mathbf{K}_1\)-双线性映射，从 \(\mathrm{M}_{(\mathbf{K}_1)} \times \mathrm{M}_{(\mathbf{K}_1)}\) 到 \(\mathrm{M}_{(\mathbf{K}_1)}\)（参见 Algebra, Chapter IX, § 1, no. 4），使 \(\mathrm{M}_{(\mathbf{K}_1)}\) 具有 \(\mathbf{K}_1\) 上的代数结构（称为由 M 将标量环扩张到 \(\mathbf{K}_1\) 后得到的代数）。特别地，当 K 是含单位元的 \(\mathbf{K}_1\) 的子环，且 \(\sigma\) 是 K 到 \(\mathbf{K}_1\) 的包含映射时，也是这种情形。

